% GENERATED FILE. Edit canonical lessons, then run npm run generate:books.
% canonical-source-hash: fnv1a64:ac3a7d81a3ee70ae
% canonical-lessons: SA-C458-darpah, SA-C458-slaghate, SA-C458-adambarah, SA-C458-uddhatah, SA-C458-vinitah

\chapter{Pride and Modesty}
\label{ch:sa-pride-and-modesty}
\hlchaptermodality{458}

\begin{chapteropening}
\textbf{By the end of this chapter:} \emph{I can say whether someone is proud or modest.}

Complete the last lesson of chapter 458: I can say whether someone is proud or modest.
\end{chapteropening}

\section[\texorpdfstring{\sk{दर्पः} (darpaḥ)}{दर्पः (darpaḥ)}]{\sk{दर्पः} (darpaḥ) — haughtiness}

\label{lesson:SA-C458-darpah}



\begin{quote}
Before the new one: say the Sanskrit for a loud laugh, then the Sanskrit for amusement.
\end{quote}

\subsection*{You'll want to know: \sk{दर्पः}}
\textbf{\sk{दर्पः}} — \emph{darpaḥ} — \textquotedblleft{}haughtiness\textquotedblright{}.

\begin{quote}
\sk{दर्पेण} \sk{सः} \sk{मित्राणि} \sk{त्यजति।}
\end{quote}

\begin{quote}
\emph{darpeṇa saḥ mitrāṇi tyajati.}
\end{quote}

\begin{quote}
In his haughtiness he drops his friends.
\end{quote}

A masculine noun.

\begin{cousinweb}
\textbf{\sk{दर्पः}} (\emph{darpaḥ}) comes from the root √dṛp, \textquotedblleft{}to be proud, to be wild\textquotedblright{}; MW gives \textquotedblleft{}pride, arrogance, haughtiness, conceit\textquotedblright{} (MW).
\end{cousinweb}

\begin{grammarlens}[title={What you've built}]
One more word for describing what you lived and felt.
\end{grammarlens}

\subsection*{Your turn}
\begin{itemize}
\raggedright
  \item \emph{Say it:} \emph{darpaḥ}
  \item \emph{Say it:} \emph{darpaḥ}, once more
  \item \emph{Say it:} \emph{darpeṇa saḥ mitrāṇi tyajati.} — In his haughtiness he drops his friends.
  \item \emph{Recall:} say the Sanskrit for a loud laugh, then the Sanskrit for amusement, then say \emph{darpaḥ} again
\end{itemize}

\begin{tcolorbox}[breakable,colback=teal!4,colframe=teal!35!black,title={Before you move on}]
What does \sk{दर्पः} mean? (\textquotedblleft{}Haughtiness\textquotedblright{}.) Say it once more.
\end{tcolorbox}
\section[\texorpdfstring{\sk{श्लाघते} (ślāghate)}{श्लाघते (ślāghate)}]{\sk{श्लाघते} (ślāghate) — boasts}

\label{lesson:SA-C458-slaghate}



\begin{quote}
Before the new one: say the Sanskrit for amusement, then the Sanskrit for haughtiness.
\end{quote}

\subsection*{You'll want to know: \sk{श्लाघते}}
\textbf{\sk{श्लाघते}} — \emph{ślāghate} — \textquotedblleft{}boasts\textquotedblright{}.

\begin{quote}
\sk{धनिकः} \sk{धनेन} \sk{श्लाघते।}
\end{quote}

\begin{quote}
\emph{dhanikaḥ dhanena ślāghate.}
\end{quote}

\begin{quote}
The rich man boasts of his wealth.
\end{quote}

Past (he …): \sk{सः} \sk{धनेन} \sk{श्लाघितवान्}, \emph{saḥ dhanena ślāghitavān}.

\begin{cousinweb}
\textbf{\sk{श्लाघते}} (\emph{ślāghate}) comes from the root √ślāgh, which MW glosses \textquotedblleft{}to talk confidently, vaunt, boast\textquotedblright{} and also \textquotedblleft{}to praise\textquotedblright{} (MW).
\end{cousinweb}

\begin{grammarlens}[title={What you've built}]
One more word for describing what you lived and felt.
\end{grammarlens}

\subsection*{Your turn}
\begin{itemize}
\raggedright
  \item \emph{Say it:} \emph{ślāghate}
  \item \emph{Say it:} \emph{ślāghate}, once more
  \item \emph{Say it:} \emph{dhanikaḥ dhanena ślāghate.} — The rich man boasts of his wealth.
  \item \emph{Recall:} say the Sanskrit for amusement, then the Sanskrit for haughtiness, then say \emph{ślāghate} again
\end{itemize}

\begin{tcolorbox}[breakable,colback=teal!4,colframe=teal!35!black,title={Before you move on}]
What does \sk{श्लाघते} mean? (\textquotedblleft{}Boasts\textquotedblright{}.) Say it once more.
\end{tcolorbox}
\section[\texorpdfstring{\sk{आडम्बरः} (āḍambaraḥ)}{आडम्बरः (āḍambaraḥ)}]{\sk{आडम्बरः} (āḍambaraḥ) — bombast}

\label{lesson:SA-C458-adambarah}



\begin{quote}
Before the new one: say the Sanskrit for haughtiness, then the Sanskrit for boasts.
\end{quote}

\subsection*{You'll want to know: \sk{आडम्बरः}}
\textbf{\sk{आडम्बरः}} — \emph{āḍambaraḥ} — \textquotedblleft{}bombast\textquotedblright{}.

\begin{quote}
\sk{तस्य} \sk{भाषणे} \sk{केवलम्} \sk{आडम्बरः} \sk{अस्ति।}
\end{quote}

\begin{quote}
\emph{tasya bhāṣaṇe kevalam āḍambaraḥ asti.}
\end{quote}

\begin{quote}
His speech is mere bombast.
\end{quote}

A masculine noun.

\begin{cousinweb}
\textbf{\sk{आडम्बरः}} (\emph{āḍambaraḥ}) MW's first sense is \textquotedblleft{}a kind of drum\textquotedblright{}, then \textquotedblleft{}a great noise\textquotedblright{} and \textquotedblleft{}noisy behaviour, bombast\textquotedblright{}; MW gives no derivation, so its origin is uncertain (MW).
\end{cousinweb}

\begin{grammarlens}[title={What you've built}]
One more word for describing what you lived and felt.
\end{grammarlens}

\subsection*{Your turn}
\begin{itemize}
\raggedright
  \item \emph{Say it:} \emph{āḍambaraḥ}
  \item \emph{Say it:} \emph{āḍambaraḥ}, once more
  \item \emph{Say it:} \emph{tasya bhāṣaṇe kevalam āḍambaraḥ asti.} — His speech is mere bombast.
  \item \emph{Recall:} say the Sanskrit for haughtiness, then the Sanskrit for boasts, then say \emph{āḍambaraḥ} again
\end{itemize}

\begin{tcolorbox}[breakable,colback=teal!4,colframe=teal!35!black,title={Before you move on}]
What does \sk{आडम्बरः} mean? (\textquotedblleft{}Bombast\textquotedblright{}.) Say it once more.
\end{tcolorbox}
\section[\texorpdfstring{\sk{उद्धतः} (uddhataḥ)}{उद्धतः (uddhataḥ)}]{\sk{उद्धतः} (uddhataḥ) — arrogant}

\label{lesson:SA-C458-uddhatah}



\begin{quote}
Before the new one: say the Sanskrit for boasts, then the Sanskrit for bombast.
\end{quote}

\subsection*{You'll want to know: \sk{उद्धतः}}
\textbf{\sk{उद्धतः}} — \emph{uddhataḥ} — \textquotedblleft{}arrogant\textquotedblright{}.

\begin{quote}
\sk{उद्धतः} \sk{छात्रः} \sk{अध्यापकं} \sk{न} \sk{नमति।}
\end{quote}

\begin{quote}
\emph{uddhataḥ chātraḥ adhyāpakaṃ na namati.}
\end{quote}

\begin{quote}
The arrogant student does not bow to the teacher.
\end{quote}

\begin{cousinweb}
\textbf{\sk{उद्धतः}} (\emph{uddhataḥ}) is built from \emph{ud}, \textquotedblleft{}up\textquotedblright{}, and the root √han, \textquotedblleft{}to strike\textquotedblright{}: raised up; MW gives \textquotedblleft{}lifted up\textquotedblright{} and then \textquotedblleft{}puffed up, haughty, arrogant, rude\textquotedblright{} (MW).
\end{cousinweb}

\begin{grammarlens}[title={What you've built}]
One more word for describing what you lived and felt.
\end{grammarlens}

\subsection*{Your turn}
\begin{itemize}
\raggedright
  \item \emph{Say it:} \emph{uddhataḥ}
  \item \emph{Say it:} \emph{uddhataḥ}, once more
  \item \emph{Say it:} \emph{uddhataḥ chātraḥ adhyāpakaṃ na namati.} — The arrogant student does not bow to the teacher.
  \item \emph{Recall:} say the Sanskrit for boasts, then the Sanskrit for bombast, then say \emph{uddhataḥ} again
\end{itemize}

\begin{tcolorbox}[breakable,colback=teal!4,colframe=teal!35!black,title={Before you move on}]
What does \sk{उद्धतः} mean? (\textquotedblleft{}Arrogant\textquotedblright{}.) Say it once more.
\end{tcolorbox}
\section[\texorpdfstring{\sk{विनीतः} (vinītaḥ)}{विनीतः (vinītaḥ)}]{\sk{विनीतः} (vinītaḥ) — modest}

\label{lesson:SA-C458-vinitah}



\begin{quote}
Before the new one: say the Sanskrit for bombast, then the Sanskrit for arrogant.
\end{quote}

\subsection*{You'll want to know: \sk{विनीतः}}
\textbf{\sk{विनीतः}} — \emph{vinītaḥ} — \textquotedblleft{}modest\textquotedblright{}.

\begin{quote}
\sk{विद्वान्} \sk{अपि} \sk{सः} \sk{विनीतः} \sk{अस्ति।}
\end{quote}

\begin{quote}
\emph{vidvān api saḥ vinītaḥ asti.}
\end{quote}

\begin{quote}
Though learned, he is modest.
\end{quote}

\begin{cousinweb}
\textbf{\sk{विनीतः}} (\emph{vinītaḥ}) is built from \emph{vi} and the root √nī, \textquotedblleft{}to lead\textquotedblright{}, the root of \textbf{\sk{नयति}}: led and trained; MW gives \textquotedblleft{}trained, well-behaved, humble, modest\textquotedblright{} (MW).
\end{cousinweb}

\begin{grammarlens}[title={What you've built}]
One more word for describing what you lived and felt.
\end{grammarlens}

\subsection*{Your turn}
\begin{itemize}
\raggedright
  \item \emph{Say it:} \emph{vinītaḥ}
  \item \emph{Say it:} \emph{vinītaḥ}, once more
  \item \emph{Say it:} \emph{vidvān api saḥ vinītaḥ asti.} — Though learned, he is modest.
  \item \emph{Recall:} say the Sanskrit for bombast, then the Sanskrit for arrogant, then say \emph{vinītaḥ} again
\end{itemize}

\begin{tcolorbox}[breakable,colback=teal!4,colframe=teal!35!black,title={Before you move on}]
What does \sk{विनीतः} mean? (\textquotedblleft{}Modest\textquotedblright{}.) Say it once more.
\end{tcolorbox}
